\begin{theorem}[General Recursion Theorem]
\label{thm:general-recursion-theorem-for-peano-system}
Let $(P,S,1)$ be a Peano system, let $W$ be a set, let $c\in W$, and let
$\Phi:P\times W\to W$ be a stage-dependent step rule. Then there exists a
unique function $f:P\to W$ such that
\[
f(1)=c
\qquad\text{and}\qquad
\forall n\in P\;(f(S(n))=\Phi(n,f(n))).
\]
\hyperref[prf:general-recursion-theorem-for-peano-system]{Go to proof.}
\end{theorem}
\end{theorembox}
\LeanFormalizes{thm:general-recursion-theorem-for-peano-system}{lra-lean}{LRA.VolumeII.PeanoSystems.Recursion.Iterator}{GeneralRecursionTheoremForPeanoSystem}{pending}

\begin{remark*}[Standard quantified statement]
\[
\exists! f:P\to W\;
\bigl(f(1)=c\land \forall n\in P\;(f(S(n))=\Phi(n,f(n)))\bigr).
\]
\end{remark*}

\begin{remark*}[Predicate reading]
\[
\exists! f:P\to W\;
\operatorname{StageDependentRecursionSolution}(f,(P,S,1),W,c,\Phi).
\]
\end{remark*}

\begin{remark*}[Negated quantified statement]
\[
\neg\exists! f:P\to W\;
\bigl(f(1)=c\land \forall n\in P\;(f(S(n))=\Phi(n,f(n)))\bigr).
\]
\end{remark*}

\begin{remark*}[Failure modes]
\begin{description}
  \item[Exposition.]
  General recursion fails if no function satisfies the general recursive
  clauses or if two distinct functions satisfy them.
  \item[Existence or uniqueness failure.]
  \[
  \neg\exists! f:P\to W\;
  \operatorname{StageDependentRecursionSolution}(f,(P,S,1),W,c,\Phi).
  \]
  \[
  \begin{aligned}
  &\neg\exists! f:P\to W\;
  \bigl(f(1)=c\land \forall n\in P\;(f(S(n))=\Phi(n,f(n)))\bigr)\\
  \Longleftrightarrow\;&
  \underbrace{\neg\exists f:P\to W\;
  \bigl(f(1)=c\land \forall n\in P\;(f(S(n))=\Phi(n,f(n)))\bigr)}_{\text{existence fails}}\\
  &\lor
  \underbrace{\exists f,h:P\to W\;
  \bigl(f\neq h\land f(1)=c\land h(1)=c
  \land \forall n\in P\;(f(S(n))=\Phi(n,f(n)))
  \land \forall n\in P\;(h(S(n))=\Phi(n,h(n)))\bigr)}_{\text{uniqueness fails}}.
  \end{aligned}
  \]
\end{description}
\end{remark*}

\begin{remark*}[Interpretation]
General recursion justifies recursive definitions whose successor rule depends
on the current stage as well as the current value. Definitions of addition,
multiplication, exponentiation, and later arithmetic operations are not part
of this section. They begin after the recursion principles have been
established.
\end{remark*}

\begin{remark*}[Historical note]
This is the stage-dependent form corresponding to Feferman's Theorem 3.4'
\citep{FefermanNumberSystems1964}. It is the form used later for arithmetic
definitions; those definitions are deferred to their own sections
\citep{FefermanNumberSystems1964,MendelsonNumberSystemsFoundationsAnalysis}.
\end{remark*}

\begin{dependencies}
\begin{itemize}
  \item \hyperref[lem:uniqueness-of-general-recursive-functions]{Uniqueness of Stage-Dependent Recursion Solutions}
  \item \hyperref[thm:general-recursion-by-state-encoding]{General Recursion by State Encoding}
\end{itemize}
\end{dependencies}

